\documentclass[12pt]{article}
\usepackage[utf8]{inputenc}
\usepackage[T1]{fontenc}
\usepackage[spanish]{babel}
\usepackage{geometry}
\usepackage{amsmath,amssymb}
\usepackage{enumitem}

\geometry{margin=2.5cm}

\newcommand{\dudoso}[1]{\textsf{#1}}

\begin{document}

\begin{center}
{\Large Transcripción cruda}\\[2mm]
{\large 14May26}
\end{center}

\section*{Foto 1}

\textit{Teorema.} Toda vecindad es un conjunto abierto.

\textit{Dem.}
Sea $p\in X$, $r>0$.
Definimos
\[
N_r(p)=\{q\in X\mid d(q,p)<r\}.
\]
Para todo $q\in N_r(p)$ existe $h>0$ tal que
\[
N_h(q)\subset N_r(p),
\]
donde se toma $h=r-d(q,p)>0$.

\section*{Foto 2}

\textit{Definici\'on.}
Sea $E\subset X$.
Se define $E'$ como el conjunto de todos los puntos l\'imite de $E$.

\textit{Teorema.}
Si $p\in E'$ entonces $N_r(p)\cap E$ es un conjunto infinito, para todo $r>0$.

\textit{Dem.}
Sup\'ongase que la afirmaci\'on es falsa. Entonces existe $r>0$ tal que $N_r(p)\cap E$ es finito.

\section*{Foto 3}

Si $p\in E'$ entonces, para cualquier $r>0$, $N_r(p)\cap E\neq\varnothing$.

\textit{Corolario.}
Si $E\subset X$ es un conjunto finito, entonces $E'=\varnothing$.

\section*{Foto 4}

\textit{Recuerdemos que un conjunto es cerrado si} $E'\subset E$.

Como $E'=\varnothing$ cuando $E$ es finito y $\varnothing\subset E$, entonces $E$ es cerrado.

\textit{Ejemplos.}
Considerar los siguientes conjuntos:
\[
A=\{(x,y)\in\mathbb R^2\mid x^2+y^2<1\}=N_1((0,0)),
\]
\[
B=\{(x,y)\in\mathbb R^2\mid x^2+y^2\le 1\},
\]
\[
C=\{a_1,a_2,\dots,a_n\}\subset\mathbb R^2,
\qquad
D=\mathbb Z.
\]

\section*{Foto 5}

\[
A=\{(x,y)\in\mathbb R^2\mid x^2+y^2<1\}=N_1((0,0)).
\]
El 0 y el 1 son puntos l\'imite de $E$.

\section*{Foto 6}

\textit{Tabla de ejemplos.}

\[
\begin{array}{c|cccc}
 & \text{cerrado} & \text{abierto} & \text{perfecto} & \text{acotado}\\
\hline
A & \text{NO} & \text{SI} & \text{NO} & \text{SI}\\
B & \text{SI} & \text{NO} & \text{SI} & \text{SI}\\
C & \text{SI} & \text{NO} & \text{NO} & \text{SI}\\
D & \text{SI} & \text{NO} & \text{NO} & \text{NO}
\end{array}
\]

\section*{Foto 7}

\textit{Teorema.}
Sea $(X,d)$ espacio m\'etrico. Sea $E\subset X$.
Entonces $E$ es abierto si y solo si $E^c$ es cerrado.

\textit{Dem.}
Sup\'ongase que $E^c$ es cerrado. P.D. que $E$ es abierto.

\section*{Foto 8}

Sea $x\in E$. Si $x$ no fuera interior, entonces para todo $r>0$ se tendr\'ia
\[
N_r(x)\cap E^c\neq\varnothing.
\]
Pero entonces $x\in (E^c)'$, contradicci\'on con que $E^c$ es cerrado.

\section*{Foto 9}

\[
N_r(x)\subset E
\]
para alg\'un $r>0$.

Rec\'iprocamente, si $E$ es abierto, entonces $E^c$ es cerrado, es decir,
\[
(E^c)'\subset E^c.
\]

\section*{Foto 10}

\textit{Teorema.} Si $n\in\mathbb N$ y $G_1,\dots,G_n\subset X$ son abiertos, entonces
\[
G_1\cap\cdots\cap G_n
\]
es abierto.

\textit{Dem.}
Sea $x\in G_1\cap\cdots\cap G_n$.
Por hip\'otesis de inducci\'on, si $G_1,\dots,G_n$ son abiertos, entonces
\[
G_1\cap\cdots\cap G_n
\]
es abierto.

\section*{Foto 11}

\textit{Teorema.} Si $n\in\mathbb N$ y $F_1,\dots,F_n\subset X$ son cerrados, entonces
\[
F_1\cup\cdots\cup F_n
\]
es cerrado.

Por inducci\'on,
\[
(F_1\cup\cdots\cup F_n)^c=F_1^c\cap\cdots\cap F_n^c
\]
es abierto.

\section*{Foto 12}

\[
\bigcap_{\alpha\in\Lambda}F_\alpha
\]
es cerrado si cada $F_\alpha$ es cerrado.

\[
\bigcup_{\alpha\in\Lambda}G_\alpha
\]
es abierto si cada $G_\alpha$ es abierto.

\section*{Foto 13}

\textit{Teorema.}
Si $F_1,\dots,F_n$ son cerrados, entonces
\[
F_1\cup\cdots\cup F_n
\]
es cerrado.

\textit{Dem.}
Por hip\'otesis de inducci\'on.

\section*{Foto 14}

\textit{Recordemos que un conjunto es cerrado si contiene a todos sus puntos l\'imite.}

Si $E$ es finito, entonces $E'=\varnothing$ y por tanto $E$ es cerrado.

\section*{Foto 15}

\textit{Ejemplos.}

\[
E=\left\{\frac1n\mid n\in\mathbb N\right\},\qquad F=\mathbb R^2,\qquad
G=(a,b),\qquad H=[a,b],\qquad I=[a,b),\qquad J=(a,b].
\]

Tabla:
\[
\begin{array}{c|cccc}
 & \text{cerrado} & \text{abierto} & \text{perfecto} & \text{acotado}\\
\hline
E & \text{NO} & \text{NO} & \text{NO} & \text{SI}\\
F & \text{SI} & \text{SI} & \text{SI} & \text{NO}\\
G & \text{NO} & \text{SI} & \text{NO} & \text{SI}\\
H & \text{SI} & \text{NO} & \text{SI} & \text{SI}\\
I & \text{NO} & \text{NO} & \text{NO} & \text{SI}\\
J & \text{NO} & \text{NO} & \text{NO} & \text{SI}
\end{array}
\]

\section*{Foto 16}

\textit{Teorema.}
Si $p\in E'$ entonces $N_r(p)\cap E$ es un conjunto infinito, para todo $r>0$.

\textit{Dem.}
Sup\'ongase que la afirmaci\'on es falsa.
Entonces existe $r>0$ tal que $N_r(p)\cap E$ es finito.

\section*{Foto 17}

Si
\[
\{q_1,\dots,q_n\}=N_r(p)\cap E,
\]
entonces tomando
\[
t:=\min\{d(q_1,p),\dots,d(q_n,p)\}>0
\]
se tiene
\[
N_t(p)\cap(E\setminus\{p\})=\varnothing,
\]
contradicci\'on.

\section*{Foto 18}

\[
N_r(p)\cap(E\setminus\{p\})\neq\varnothing.
\]
\[
E'\subset E
\]
si $E$ es cerrado.

\section*{Foto 19}

\textit{Corolario.} Si $E\subset X$ es un conjunto finito, entonces $E'=\varnothing$.

\textit{Teorema.} Un conjunto $E\subset X$ es perfecto si $E=E'$ y es cerrado.

\section*{Foto 20}

\textit{Conclusi\'on.}
Si $E$ es abierto entonces $E^c$ es cerrado.
Si $E^c$ es cerrado entonces $E$ es abierto.

\end{document}
